% ============================================================
%  Chapitre 5 — Logique Classique
% ============================================================
\chapter{Logique Classique}

\section{Logique propositionnelle}

\subsection{Fondements}

La \textbf{logique propositionnelle} est le formalisme le plus élémentaire.
Une formule est construite à partir d'atomes $p, q, r, \ldots$ et des
connecteurs $\lnot, \land, \lor, \to, \leftrightarrow$.

\noindent La \textbf{forme normale conjonctive} (CNF) est une conjonction de
clauses, chaque clause étant une disjonction de littéraux :
\[
  \underbrace{(l_{11} \lor \ldots \lor l_{1k_1})}_{\text{clause }C_1}
  \land \ldots \land
  \underbrace{(l_{m1} \lor \ldots \lor l_{mk_m})}_{\text{clause }C_m}
\]

\noindent Le problème de \textbf{satisfiabilité} (\textsc{sat}) — trouver une
valuation qui rend la formule vraie — est NP-complet.
Le solveur \textbf{SAT4J} résout ce problème de manière efficace en pratique.

\subsection{Application — Industrie musicale}

Les formules modélisent des contraintes sur les artistes et les albums :

\begin{center}
\begin{tabular}{ll}
\toprule
\textbf{Atome} & \textbf{Signification} \\
\midrule
\code{SigneLabel}       & L'artiste est signé chez un label \\
\code{LibreDroits}      & L'artiste est libre de droits \\
\code{AlbumSorti}       & L'artiste a sorti un album \\
\code{Royalties}        & L'artiste perçoit des royalties \\
\code{AlbumCertifie}    & L'album est certifié \\
\bottomrule
\end{tabular}
\end{center}

\noindent\textbf{Clauses :}
\begin{align*}
C_1 &= \lnot\text{SigneLabel} \lor \lnot\text{LibreDroits}
    && \text{(signé $\Rightarrow$ pas libre de droits)} \\
C_2 &= \lnot\text{AlbumSorti} \lor \text{Royalties}
    && \text{(album sorti $\Rightarrow$ royalties)} \\
C_3 &= \lnot\text{SigneLabel} \lor \text{Royalties}
    && \text{(signé $\Rightarrow$ royalties)} \\
C_4 &= \text{SigneLabel} \lor \text{LibreDroits}
    && \text{(signé ou libre)} \\
C_5 &= \lnot\text{AlbumCertifie} \lor \text{AlbumSorti}
    && \text{(certifié $\Rightarrow$ sorti)}
\end{align*}

\subsection{Application — Système judiciaire}

\begin{align*}
C_1 &= \lnot\text{Coupable} \lor \lnot\text{Innocent}
    && \text{(exclusion mutuelle)} \\
C_2 &= \text{Coupable} \lor \text{Innocent}
    && \text{(tiers exclu)} \\
C_3 &= \lnot\text{Recidiviste} \lor \text{PeineLourde}
    && \text{(récidiviste $\Rightarrow$ peine lourde)} \\
C_4 &= \lnot\text{Mineur} \lor \lnot\text{CourAssises}
    && \text{(mineur pas en cour d'assises)} \\
C_5 &= \lnot\text{Temoignage} \lor \text{AcquitPossible}
    && \text{(témoignage $\Rightarrow$ acquittement possible)}
\end{align*}

\subsection{Implémentation avec SAT4J}

\begin{lstlisting}[caption={Résolution SAT propositionnelle via Tweety}]
PlParser parser = new PlParser();
PlBeliefSet kb = new PlBeliefSet();

// Ajout des clauses (formules)
kb.add(parser.parseFormula("!SigneLabel || !LibreDroits"));
kb.add(parser.parseFormula("!AlbumSorti || Royalties"));
kb.add(parser.parseFormula("SigneLabel || LibreDroits"));

// Requete de satisfiabilite
Sat4jSolver solver = new Sat4jSolver();
boolean sat = solver.isSatisfiable(kb);
// → true : la base est cohérente

// Inference : LibreDroits est-il entraine par SigneLabel=faux ?
PlFormula query = parser.parseFormula("LibreDroits");
boolean entailed = solver.isEntailed(kb, query);
\end{lstlisting}

\section{Logique des prédicats du premier ordre}

\subsection{Fondements}

La \textbf{logique des prédicats du premier ordre} (\textsc{fol}) enrichit la
logique propositionnelle avec des variables, des prédicats, des fonctions et des
quantificateurs $\forall$ (universel) et $\exists$ (existentiel).

\noindent Le \textbf{chaînage avant} (\emph{forward chaining}) est une stratégie
d'inférence qui part des faits connus et applique les règles disponibles pour
dériver de nouveaux faits, jusqu'à stabilisation.

\subsection{Application — Industrie musicale}

\textbf{Faits initiaux :}
\begin{align*}
&\text{Artiste}(\text{soolking}),\quad \text{SigneLabel}(\text{soolking}),\quad
 \text{AlbumSorti}(\text{soolking}) \\
&\text{Artiste}(\text{tif}),\quad \text{AlbumSorti}(\text{tif})
\end{align*}

\textbf{Règles :}
\begin{align*}
R_1&: \forall x,\; \text{SigneLabel}(x) \Rightarrow \lnot\text{LibreDroits}(x) \\
R_2&: \forall x,\; \text{Artiste}(x) \land \text{AlbumSorti}(x)
       \Rightarrow \text{Royalties}(x) \\
R_3&: \forall x,\; \text{SigneLabel}(x) \Rightarrow \text{Royalties}(x)
\end{align*}

\textbf{Dérivations (chaînage avant) :}
\begin{align*}
R_1 + \text{SigneLabel(soolking)} &\Rightarrow \lnot\text{LibreDroits(soolking)} \\
R_2 + \text{Artiste(soolking)} + \text{AlbumSorti(soolking)} &\Rightarrow \text{Royalties(soolking)} \\
R_2 + \text{Artiste(tif)} + \text{AlbumSorti(tif)} &\Rightarrow \text{Royalties(tif)}
\end{align*}

\subsection{Application — Système judiciaire}

\textbf{Faits initiaux :}
\[
  \text{Accuse(ali)},\quad \text{Accuse(omar)},\quad \text{Recidiviste(omar)},\quad
  \text{Mineur(karim)},\quad \text{Temoignage(ali)}
\]

\textbf{Règles :}
\begin{align*}
R_1&: \forall x,\; \text{Recidiviste}(x) \Rightarrow \text{PeineLourde}(x) \\
R_2&: \forall x,\; \text{Mineur}(x) \Rightarrow \text{TribunalEnfants}(x) \\
R_3&: \forall x,\; \text{Temoignage}(x) \Rightarrow \text{AcquitPossible}(x)
\end{align*}

\subsection{Implémentation}

\begin{lstlisting}[caption={Chaînage avant avec Tweety FOL}]
FolParser parser = new FolParser();
FolBeliefSet kb = new FolBeliefSet();

// Predicats
kb.getSignature().add(new Predicate("Artiste", 1));
kb.getSignature().add(new Predicate("SigneLabel", 1));
kb.getSignature().add(new Predicate("Royalties", 1));

// Faits
kb.add(parser.parseFormula("Artiste(soolking)"));
kb.add(parser.parseFormula("SigneLabel(soolking)"));

// Regles (implication universelle)
kb.add(parser.parseFormula("forall X: (SigneLabel(X) => Royalties(X))"));

// Chainage avant
NaivePrologReasoner reasoner = new NaivePrologReasoner();
Collection<FolFormula> derived = reasoner.query(kb,
    parser.parseFormula("Royalties(soolking)"));
\end{lstlisting}
